\section{Define the encoded instructions}\label{vm-opcodes}
Opcode bytes and encoded byte sequences in this chapter are hexadecimal. Byte offsets,
lengths, depths and instruction counts are decimal unless a prefix says otherwise.
There are thirteen legal opcodes. Each instruction consists of one opcode followed
by its fixed number of operand bytes.

\begin{table}[htbp]
\centering
\caption{The exercise's guest instruction set. Stack effects use the earlier live value as the left arithmetic operand.}
\tagpdfsetup{table/header-rows={1}}
\begin{tabular}{llrp{80mm}}
\toprule
Opcode & Name & Bytes & Effect\\
\midrule
\texttt{00} & \texttt{HALT} & 1 & Stop successfully; retain live operands.\\
\texttt{01} & \texttt{PUSH8} & 2 & Push the following signed byte.\\
\texttt{02} & \texttt{PUSH32} & 5 & Push the following signed little-endian word.\\
\texttt{10} & \texttt{ADD} & 1 & Replace left and right with their checked sum.\\
\texttt{11} & \texttt{SUB} & 1 & Replace them with checked left minus right.\\
\texttt{12} & \texttt{MUL} & 1 & Replace them with their checked product.\\
\texttt{13} & \texttt{DIV} & 1 & Replace them with the checked signed quotient.\\
\texttt{20} & \texttt{DUP} & 1 & Copy the top live value onto the stack.\\
\texttt{21} & \texttt{DROP} & 1 & Remove the top live value.\\
\texttt{22} & \texttt{SWAP} & 1 & Exchange the two top live values.\\
\texttt{23} & \texttt{OVER} & 1 & Copy the second live value from the top.\\
\texttt{30} & \texttt{JMP} & 2 & Branch to the following unsigned byte offset.\\
\texttt{31} & \texttt{JZ} & 2 & Pop a test value; branch if it is zero.\\
\bottomrule
\end{tabular}
\end{table}

A branch operand is an absolute offset from the program base. The byte after
\texttt{JMP} is neither a signed displacement nor an address. \texttt{JZ} consumes
its test value on both outcomes: zero selects the encoded target, while a nonzero
value selects the byte after the two-byte branch instruction.

\texttt{PUSH8} sign extends its operand: \texttt{80} pushes minus 128 and
\texttt{ff} pushes minus one. \texttt{PUSH32} reads four bytes in least-significant
first order. Bytes \texttt{00 00 00 80} encode the minimum signed word directly.
All 32-bit patterns are valid constants. This format has no decimal-literal parser
or separate unary-sign token. The image's starting address need not be word aligned,
so the guided interpreter assembles that word from individual bytes.
